\chapter{一般群的 Langlands 纲领 II}

% ===== PDF p.094 =====
\noindent\textbf{上次}
\begin{itemize}
\item $F$ 为任意数域。
\item $\GL_{n,F}$ 的 Langlands（–Clozel–Fontaine–Mazur）互反猜想断言：
\begin{itemize}
\item $\ell$ 为某个素数。固定一个同构 $\C\cong\Ql$。
\item 则对任意 $n\ge1$，存在一个（唯一的）双射，介于
\begin{itemize}
\item[(i)] $\GL_n(\A_F)$ 的 $L$-代数尖点自守表示之集，
\item[(ii)] 不可约连续表示 $\Gal(\overline F/F)\to\GL_n(\Ql)$ 的同构类之集，这些表示几乎处处非分歧，且在整除 $\ell$ 的位处为 de Rham 表示，
\end{itemize}
之间，使得该双射把 Satake 参数\idx{Satake 参数}与 Frobenius 元素的特征值相匹配。
\end{itemize}
\item 对任意 $G$ 引入了（阿代尔）自守形式的空间 $\cA(G)$。
\item 由 Flath 定理 $\Rightarrow$ 自守表示 $\pi$ 分解：$\pi\cong\sideset{}{'}\bigotimes_p\pi_p\otimes\pi_\infty$。
\item 解释了「$L$-代数」（= $\pi_\infty$ 的无穷小特征是「整的」）。
\end{itemize}

\medskip
\noindent\textbf{今日}
\begin{itemize}
\item 引入「尖点」自守表示。
\item 解释 Satake 参数。
\item 讨论一般 $G$ 的 Langlands 纲领的一些预期。
\end{itemize}

\medskip
\noindent\textbf{定义 12.1}（\texttt{[GH19, Definition 9.2]}）\phantomsection\label{def:12.1}\quad 设 $G/\Q$ 约化。称 $\phi\in L^2([G])$（相应地 $\phi\in\cA(G)$）为\textbf{尖点的}，若对一切真抛物子群 $P\subseteq G$（定义在 $\Q$ 上）的幺幂根基 $N$ 以及几乎所有 $g\in G(\A)$，有
\begin{equation}\tag{1}\label{eq:12.1}
\int_{N(\Q)\backslash N(\A)}\phi(ng)\,\mathrm{d}n=0.
\end{equation}

\medskip
\noindent\textbf{尖点自守表示}
\begin{itemize}
\item 记 $L^2_{\mathrm{cusp}}([G])$、相应地 $\cA_{\mathrm{cusp}}(G)$、$\cA_{\mathrm{cusp}}([G])$，为尖点元素构成的子空间。
\item 存在像稠密的嵌入
\[
\cA_{\mathrm{cusp}}([G])\subseteq L^2_{\mathrm{cusp}}([G]),
\]
即：尖点自守形式满足足够强的增长条件，使其属于 $L^2$（它们是「速降」的），参见 \texttt{[GH19, Section 6.5.]}。
\item $G(\A)$ 的 $L^2$-自守表示称为\textbf{尖点的}，若它同构于 $L^2_{\mathrm{cusp}}([G])$ 的某个子商（实际上是子表示）。
\item 类似地，自守表示称为尖点的，若它作为 $\cA_{\mathrm{cusp}}(G)$ 的子商出现。
\item Gelfand、Piatetski-Shapiro：作为酉 $G(\A)$-表示，
% ===== PDF p.095 =====
\begin{equation}\tag{2}\label{eq:12.2}
L^2_{\mathrm{cusp}}([G])\cong\widehat{\bigoplus}_{\pi\in\widehat{G(\A)}}m_\pi\,\pi
\end{equation}
其中每个 $m_\pi$ 有限，即 $m_\pi\in\N\cup\{0\}$，参见 \texttt{[GH19, Corollary 9.1.2.]}。
\item 特别地，$L^2_{\mathrm{cusp}}([G])\subset L^2_{\mathrm{disc}}([G])$（注意：平凡表示出现在 $L^2_{\mathrm{disc}}([G])\setminus L^2_{\mathrm{cusp}}([G])$ 中）。
\item 由
\[
L^2_{\mathrm{cusp}}([G])\cong\widehat{\bigoplus}_{\pi\in\widehat{G(\A)}}m_\pi\,\pi
\]
推出分解（分解为不可约容许 $G(\Af)\times(\g_\C,K_\infty)$-模）
\begin{equation}\tag{3}\label{eq:12.3}
\cA_{\mathrm{cusp}}([G])\cong\bigoplus_{\pi\in\widehat{G(\A)}}m_\pi\,\pi^{K_\infty\text{-}\mathrm{finite}},
\end{equation}
其中 $\pi^{K_\infty\text{-}\mathrm{finite}}\subseteq\pi$ 记 $K_\infty$-有限（特别是光滑）向量构成的（稠密）子空间，参见 \texttt{[GH19, Theorem 4.4.4.]}。
\item 于是，$L^2_{\mathrm{cusp}}([G])$ 的分解与 $\cA_{\mathrm{cusp}}([G])$ 的分解「相同」。
\item Piatetski-Shapiro、Shalika，参见 \texttt{[GH19, Theorem 11.3.4]}：设 $F$ 为数域、$G=\Res_{F/\Q}\GL_{n,F}$。则 $L^2_{\mathrm{cusp}}([G])$ 上的 $G(\A)$-表示是重数自由的（「重数一」）。
\item 「强重数一」：设 $\pi,\pi'$ 是 $\GL_{n,F}$ 的尖点自守表示。若对几乎所有素数 $p$ 有 $\pi_p\cong\pi'_p$，则 $\pi\cong\pi'$，参见 \texttt{[GH19, 11.7.2]}。
\item 去掉尖点性后不成立，参见 \texttt{[BG10, page 38]}。
\item Mœglin、Waldspurger，参见 \texttt{[GH19, Theorem 10.7.1.]}、\texttt{[MW89]}：设 $G=\Res_{F/\Q}\GL_{n,F}$。则在把 $d\mid n$ 的 $\GL_{d,F}$ 的尖点自守表示参数化的意义下，$L^2_{\mathrm{disc}}([G])$ 是已知的，并且它被证明是重数自由的。
\end{itemize}

我们现在可以补上一个缺口。即：我们从未给新形式配上自守表示。但用现在陈述的材料，我们至少可以给出（两种）形式定义。设 $f\in S_k$ 为新形式。则我们可以给 $f$ 配上 $\cA(\GL_2)$ 中由自守形式 $\phi_f$ 生成的（不可约）$\GL_2(\Af)\times(\mathfrak{gl}_{2,\C},\mathrm{O}_2(\R))$-子模（参见第 3 节），也可以配上 $L^2([\GL_2])$ 中由 $\tilde\phi_f$ 生成的不可约酉子表示（参见第 3 节）。具体用哪种构造，要在相应的具体情形中检验。

\medskip
\noindent\textbf{Satake 参数}
\begin{itemize}
\item 描述 $\pi$ 在 $p$ 处非分歧时的 $\pi_p$，即 $\pi_p^{G_p(\Z_p)}\ne0$，其中 $G_p\to\Spec(\Z_p)$ 是 $G_{\Q_p}$ 的约化模型（对几乎所有素数 $p$ 存在）。
\item 也就是描述 $\C$-代数同态
\[
\cH(G(\Q_p),K)\longrightarrow\C,\qquad K:=G_p(\Z_p).
\]
\item Langlands/Satake，参见 \texttt{[GH19, Theorem 7.5.1., Corollary 7.5.2.]}、\texttt{[BG10, Section 2.1.]}：存在
\begin{itemize}
\item 一个「自然的」$\Q$ 上代数群 $\widehat G$，
\item 一个自同构 $\Frob_p\colon\widehat G(\Q)\cong\widehat G(\Q)$，
% ===== PDF p.096 =====
\item 以及一个双射
\begin{equation}\tag{4}\label{eq:12.4}
\Hom_{\C\text{-}\mathrm{alg}}\bigl(\cH(G(\Q_p),K),\C\bigr)\;\overset{1:1}{\longleftrightarrow}\;\bigl\{\widehat G(\C)\rtimes\Frob_p^{\Z}\ \text{中的 Frobenius 半单}\ \widehat G(\C)\text{-共轭类}\bigr\}.
\end{equation}
\end{itemize}
\item 群 $\widehat G_\C$ 与 $G$ 具有对偶的根数据。因此例如（参见 \texttt{[Bor79]}）：
\begin{itemize}
\item 若 $G=\GL_n$，则 $\widehat G=\GL_{n,\C}$；
\item 若 $G=\SL_n$，则 $\widehat G=\mathrm{PGL}_{n,\C}$；
\item 若 $G=\SO_{2n}$，则 $\widehat G=\SO_{2n,\C}$；
\item 若 $G=\SO_{2n+1}$，则 $\widehat G=\mathrm{Sp}_{2n,\C}$；
\item 若 $G=\mathrm{GSp}_{2n}$，则 $\widehat G=\mathrm{GSpin}_{2n+1,\C}$。
\end{itemize}
\item 若 $G_{\Q_p}$ 分裂，则自同构 $\Frob_p$ 平凡，而 Frobenius 半单共轭类就是 $\widehat G(\C)$ 中半单元素的共轭类。
\item 若 $G=\GL_n$，则半单共轭类由其元素的特征多项式唯一确定。
\item 设 $\pi=\sideset{}{'}\bigotimes_p\pi_p\otimes\pi_\infty$ 是 $G$ 的自守表示。则我们得到下述「Hecke 特征系统」的类比：
\begin{itemize}
\item 素数的一个有限集 $S$，使得 $p\notin S$ 时 $\pi_p$ 非分歧，
\item 对每个 $p\notin S$，$\widehat G(\C)\rtimes\Frob_p^{\Z}$ 中一个 Frobenius 半单共轭类 $c_p(\pi)$。
\end{itemize}
\item 对 $\GL_n$，$c_p(\pi)$ 中（每个元素的）特征值就是 $\pi$ 的 Satake 参数。
\item 具体地，若双陪集
\[
\GL_n(\Z_p)\diag\bigl(\underbrace{p,\dots,p}_{i\ \text{个}},\underbrace{1,\dots,1}_{(n-i)\ \text{个}}\bigr)\GL_n(\Z_p)
\]
在 $\pi_p^{\GL_n(\Z_p)}$ 上的特征值为 $e_{a_{p,i}}$，则 $c_p(\pi)$ 中元素的特征多项式为
\begin{equation}\tag{5}\label{eq:12.5}
X^n-p^{\frac{1-n}{2}}e_{a_{p,1}}X^{n-1}+\cdots+(-1)^ip^{\frac{i(i-n)}{2}}e_{a_{p,i}}X^i+\cdots+(-1)^ne_{a_{p,n}},
\end{equation}
参见 \texttt{[GH19, Section 7.2.]}。
\item 更具体地，对由新形式 $f=\sum_{n\ge1}a_nq^n\in S_k(\Gamma_0(N),\chi)$ 生成的 $\pi$，$p\nmid N$ 时的 $c_p(\pi)$ 具有特征多项式
\begin{equation}\tag{6}\label{eq:12.6}
X^2-p^{-1/2}p\,a_pX+\chi(p)p^k.
\end{equation}
\end{itemize}

回忆 $\pi\otimes_\C|\det|^{1/2}_{\text{adèlic}}$ 是 $L$-代数的。对 $\pi\otimes_\C|\det|^{1/2}_{\text{adèlic}}$ 得到多项式
\[
X^2-a_pX+\chi(p)p^{k-1},
\]
（在选定同构 $\Ql\cong\C$ 之后）它曾是 $p$ 处算术 Frobenius 的特征多项式。\znote{原文第 (6) 式确作 $X^2-p^{-1/2}p\,a_pX+\chi(p)p^k$，即一次项系数为 $p^{1/2}a_p$，而非 $p^{-1/2}a_p$。按此处的 Satake 参数归一化，扭 $|\det|^{1/2}$ 会把每个 Satake 参数乘上 $p^{-1/2}$：于是 $p^{1/2}a_p\cdot p^{-1/2}=a_p$，$\chi(p)p^k\cdot p^{-1}=\chi(p)p^{k-1}$，正得下一式。两式自洽，故此处并非笔误。[来源：底本 p.~96 原式]}

\medskip
\noindent\textbf{L-群}
\begin{itemize}
\item $G$ 为 $\Q$ 上的约化群。
\item $\ell$ 为素数。
\item 稍作工作（参见 \texttt{[GH19, Section 7.3.]}）可知，群 $\Gal(\overline\Q/\Q)$ 作用在 $\Q$ 上的约化群 $\widehat G$ 上（若 $G$ 分裂，则此作用平凡）。
% ===== PDF p.097 =====
\item 置
\begin{equation}\tag{7}\label{eq:12.7}
{}^LG:=\widehat G(\Ql)\rtimes\Gal(\overline\Q/\Q).
\end{equation}
\item 所谓\textbf{L-参数}，按定义是连续同态
\[
\Gal(\overline\Q/\Q)\longrightarrow{}^LG
\]
的 $\widehat G(\Ql)$-共轭类，这些同态到 $\Gal(\overline\Q/\Q)$ 的投影为恒等映射。
\item 若 $F/\Q$ 有限且 $G=\Res_{F/\Q}\GL_{n,F}$，则一个 L-参数等同于 $\Gal(\overline F/F)$ 的一个 $n$ 维 $\ell$-进 Galois 表示的同构类。
\end{itemize}

\medskip
\noindent\textbf{关于 $L$-代数自守表示的 Buzzard–Gee 猜想}
\begin{itemize}
\item 设 $\ell$ 为素数。
\item 固定一个同构 $\iota\colon\C\cong\Ql$。
\item 设 $\pi$ 是 $G$ 的 $L$-代数自守表示。
\item 则 Buzzard–Gee（参见 \texttt{[BG10, Conjecture 3.2.2.]}）猜想存在 L-参数
\[
\rho_\pi\colon\Gal(\overline\Q/\Q)\longrightarrow{}^LG,
\]
使得（特别地）
\begin{itemize}
\item 若 $p\ne\ell$ 且 $\pi$ 在 $p$ 处非分歧，则每个算术 Frobenius $\Frob_p\in\Gal(\overline{\Q}_p/\Q_p)$ 在
\[
\Gal(\overline{\Q}_p/\Q_p)\longrightarrow\widehat G(\Ql)\rtimes\Gal(\overline{\Q}_p/\Q_p)\twoheadrightarrow\widehat G(\Ql)\rtimes\Frob_p^{\widehat\Z}
\]
下的像落在 $\iota(c_p(\pi))\in\widehat G(\Ql)\rtimes\Frob_p^\Z$ 的共轭类中；
\item 对每个连续表示 ${}^LG\to\GL_N(\Ql)$（它在 $\widehat G(\Ql)$ 上是代数的），复合
\[
\Gal(\overline{\Q}_\ell/\Q_\ell)\xrightarrow{\ \rho_\pi\ }{}^LG\longrightarrow\GL_N(\Ql)
\]
是 de Rham 的。
\end{itemize}
\item 这推广了「把 Galois 表示附着到某个新形式所对应的（某个扭的）自守表示上」这一构造。
\item 与 $\GL_n$ 的情形相反，这里并不猜想 L-参数 $\rho_\pi$ 唯一，参见 \texttt{[BG10, Remark 3.2.4.]}。
\item 此外，对不同的 $L$-代数自守表示 $\pi$，我们可能得到同一个 L-参数 $\rho$。
\item 对 $\GL_n$，Buzzard–Gee 猜想在此前的 Langlands–Clozel–Fontaine–Mazur 猜想之外，还预言存在附着于（可能非尖点的）$L$-代数自守表示的 Galois 表示。此时这些 Galois 表示不再被猜想为不可约。
\end{itemize}

\medskip
\noindent\textbf{Langlands 群\idx{Langlands 群} $\cL$}
\begin{itemize}
\item 猜想存在一个（非常大的）局部紧群 $\cL$，即（整体）\textbf{Langlands 群}，具有（至少）下述性质，参见 \texttt{[GH19, Section 12.6.]}、\texttt{[Art02]}、\texttt{[LR87]}：\znoteplain{原文把这个群记作花体的 $L$，与 $L$-函数、L-群等术语中的字母 $L$ 同形，容易混淆；译文一律记作 $\cL$。}
\begin{itemize}
% ===== PDF p.098 =====
\item 存在典范满射 $\cL\twoheadrightarrow\Gal(\overline\Q/\Q)$（事实上 $\cL\twoheadrightarrow W_\Q$，其中 $W_\Q$ 是整体 Weil 群）。
\item 对每个 $n\ge0$，$\cL$ 的连续不可约 $n$ 维 $\C$-表示的等类 $\mathrm{Irr}_n(\cL)$ 自然地与 $\GL_n(\A)$ 的尖点自守表示之集 $\mathrm{Csp}_n$ 一一对应。
\item 固定一个同构 $\Ql\cong\C$。则存在单射，从不可约、几乎处处非分歧且在 $\ell$ 处为 de Rham 的 $\ell$-进 Galois 表示 $\Gal(\overline\Q/\Q)$ 的同构类之集 $\mathcal{G}_{\mathrm{geom}}$，映入 $\cL$ 的不可约表示的等类之集 $\mathrm{Irr}$，其像对应于 $L$-代数尖点自守表示。
\item 对 $\Q$ 的每个位 $v$，存在局部 Langlands 群的单射 $\cL_v\to\cL$，它在 $\cL$ 中共轭意义下良定义。这里：
\begin{itemize}
\item 若 $v=\infty$，则 $\cL_v\cong W_\R$，即 $\C^\times$ 被 $\Gal(\C/\R)$ 的非分裂扩张；
\item 若 $v=p$ 为素数，则 $\cL_v\cong W_{\Q_p}\times\mathrm{SU}_2(\R)$，其中
\[
W_{\Q_p}\cong I_{\Q_p}\rtimes\Z\subseteq\Gal(\overline{\Q}_p/\Q_p)\cong I_{\Q_p}\rtimes\widehat\Z
\]
是映入 Frobenius 的整数幂的元素构成的子群。
\end{itemize}
\item 若 $\pi\cong\sideset{}{'}\bigotimes_v\pi_v\in\mathrm{Csp}_n$ 对应于连续表示 $\rho_\pi\colon\cL\to\GL_n(\C)$，则在 $\GL_n$ 的局部 Langlands 对应下，$\pi_v$ 应当对应于 $\rho_\pi|_{\cL_v}$，参见 \texttt{[GH19, Section 12.5.]}。
\end{itemize}
\item 若这些条件成立，则必然有 $\cL^{\mathrm{ab}}\cong\Q^\times\backslash\A^\times$。注意由整体类域论存在满射 $\Q^\times\backslash\A^\times\twoheadrightarrow\Gal(\overline\Q/\Q)^{\mathrm{ab}}$。
\item $\cL$ 的存在性目前完全无法企及；更进一步地，猜想 $\cL$ 应满射到 $\C$ 上的「动机 Galois 群」$G_{\mathrm{mot}}$，参见 \texttt{[Art02]}、\texttt{[LR87]}。
\end{itemize}

\medskip
\noindent\textbf{一般 $G$ 的整体 Langlands 对应（非常、非常粗略）}
\begin{itemize}
\item 这里不追求任何精确性，所有陈述都近乎空泛！
\item 精确的陈述可查阅 \texttt{[GH19, Chapter 12.]} 或 \texttt{[Art94]}。
\item 为了参数化非 $L$-代数的自守表示（至少是那些与 $L^2_{\mathrm{disc}}([G])$ 的分解有关的），应当把此前 Galois 版本的 L-参数换成（某些）L-参数
\[
\cL\longrightarrow{}^LG
\]
（实际上是 L-群的 Weil 形式 $W_\Q\ltimes\widehat G(\C)$）。
\item 于是人们希望构造满射（参见 \texttt{[GH19, Conjecture 12.6.2.]}）
\[
\{(\text{某些})\ \text{自守表示}\}\ \xrightarrow{\ \mathcal{LL}\ }\ \{(\text{某些})\ \text{L-参数}\}.
\]
\item $\mathcal{LL}$ 的纤维称为\textbf{L-包}（L-packet），它们可能无限。
\item 局部 Langlands 对应的精确构造（参见 \texttt{[GH19, Conjecture 12.5.1.]}）应当给出 L-包的参数化，参见 \texttt{[GH19, (12.23)]}。
% ===== PDF p.099 =====
\item 对每个 $\pi\subseteq L^2_{\mathrm{disc}}([G])$，其重数应当可以通过 L-参数 $\mathcal{LL}(\pi)$ 与 L-包的精确参数化算出，参见 \texttt{[GH19, Conjecture 12.6.3.]}。
\item 然而，对非缓增（non-tempered）表示，上述做法应当不合理；此时应考虑 Arthur 参数
\[
\cL\times\SL_2(\C)\longrightarrow{}^LG
\]
代替 L-参数，参见 \texttt{[Art94]}。
\end{itemize}

\medskip
\noindent\textbf{本讲未提及的重要内容}
\begin{itemize}
\item L-函数，\texttt{[Bor79]}、\texttt{[GH19, Chapter 11, Section 12.7.]}；
\item 函子性，\texttt{[GH19, Section 12.6.]}；
\item 迹公式，\texttt{[Art05]}、\texttt{[GH19, Section 18]}；
\item 函数域的 Langlands 纲领；
\item ……
\end{itemize}

\medskip
\noindent\textbf{Shimura 簇\idx{Shimura 簇}一瞥}
\begin{itemize}
\item 设 $K_\infty\subseteq G(\R)$ 为极大的（连通）紧子群。
\item 设 $K\subseteq G(\Af)$ 为（充分小的）紧开拓子群。
\item 回忆：
\begin{itemize}
\item $[G]\to[G]/K=G(\Q)A_G\backslash G(\A)/K$ 是实流形上的射有限覆盖；
\item $[G]/K\to X_K:=[G]/KK_\infty=G(\Q)A_G\backslash G(\A)/KK_\infty$ 是 $X_K$ 上的 $K_\infty$-丛，而 $X_K$ 是算术流形的无交并（= 对称空间对算术子群的商）。
\end{itemize}
\item 注意 $X_K$ 隐含地依赖于 $K_\infty$。
\item 为简化后面的一些记号，置
\[
\widehat X:=\varprojlim_{K}X_K.
\]
\end{itemize}

\medskip
\noindent\textbf{对 $G=\GL_2$ 已出现过的东西}
\begin{itemize}
\item[(1)] 上半/下半平面 $\HH^\pm\cong G(\R)/K_\infty$，它自然地是复流形；
\item[(2)] 全纯嵌入 $\HH^\pm\to\Pone_\C$；
\item[(3)] $\HH^\pm$ 上 $G(\R)$-等变的向量丛 $\omega_k$，或经拉回成为 $K\subseteq G(\Af)$ 紧开拓时 $X_K$ 上的向量丛；
\item[(4)] 对 $k\in\Z$，模形式空间 $M_k\subseteq H^0(\widehat X,\omega_k)$，其中右端指全纯截面；
\item[(5)] $X_K$ 的典范紧化 $X_K^*$，以及 $\omega_k$ 在其上的延拓；
\item[(6)] 一个概形 $\widehat X\to\Spec(\Q)$，其 $\C$-值点自然地同构于 $\widehat X$；
% ===== PDF p.100 =====
\item[(7)] 由 $\widetilde X$ 上某些 $G(\Af)$-等变的 $\Ql$-局部系统 $\mathbb L$（例如 $\mathbb L\cong\Ql$）得到的 $\Gal(\overline\Q/\Q)\times G(\Af)$-表示\znote{原文在第 (4)、(6) 条用 $\widehat X$，在第 (7)–(9) 条改用 $\widetilde X$，所指应当相同（即前述 $\varprojlim_KX_K$）；译文保留原文的两种写法。}
\[
\widetilde H^1_{\mathrm{\acute{e}t}}\bigl(\widetilde X_{\overline\Q},\mathbb L\bigr);
\]
\item[(8)] Eichler–Shimura 同构：把 $k\ge2$ 的 $S_k$ 与某些 $\widetilde H^1_{\mathrm{\acute{e}t}}(\widetilde X_{\overline\Q},\mathbb L)$ 联系起来；
\item[(9)] Eichler–Shimura 关系：表达 $\widetilde H^1_{\mathrm{\acute{e}t}}(\widetilde X_{\overline\Q},\mathbb L)$ 上 $\Gal(\overline\Q/\Q)$-作用与 $G(\Af)$-作用之间的关系，它是某条模曲线的 mod $p$ 几何的推论；
\item[(10)] 一个映射
\[
\mathcal{LL}\colon A_{\mathrm{mod}}\longrightarrow G_{\mathrm{mod}},
\]
其像（猜想）可以描述。
\end{itemize}

对一般的 $G$，可惜并没有类比——但是等等：存在能给出 Shimura 簇的 $G$！

\begin{itemize}
\item 甚至第 (1) 条对一般的 $G$ 就失效：空间 $G(\R)/K_\infty$ 不必是复流形！（例如 $G=\GL_3$ 时，实维数
\[
\dim\bigl(\GL_3(\R)/A_{\GL_3}(\R)\SO_3(\R)\bigr)=\frac{3(3+1)}{2}-1=5
\]
是奇数。）
\item 特别地，一切与全纯性有关的东西对这样的 $G$ 都没有类比，例如第 (2)、(3)、(4)、(6)、(7) 条等。
\item 这是坏消息！
\item 好消息：对能给出 Shimura 数据\idx{Shimura 数据}的群 $G$（例如 $F/\Q$ 全实时的 $\Res_{F/\Q}\GL_{2,F}$，或 $\mathrm{GSp}_{2n}$ 等，参见 \texttt{[Lan17]}、\texttt{[Mil05]}、\texttt{[Del71a]}、\texttt{[Del79]}），实流形
\[
G(\R)/K_\infty A_G
\]
是复流形。
\item 一个 \textbf{Shimura 数据}是一对 $(G,X)$，其中 $G$ 是 $\Q$ 上的约化群，$X$ 是态射
\[
h\colon\mathbb S:=\Res_{\C/\R}\Gm\longrightarrow G_\R
\]
的 $G(\R)$-共轭类，使得
\begin{itemize}
\item[1)] 对每个 $h\in X$，$\mathbb S(\R)=\C^\times$ 经 $\mathrm{Ad}\circ h$ 在 $\Lie(G_\R)_\C$ 上作用的特征标或为 $z\mapsto z/\bar z$，或为 $z\mapsto1$，或为 $z\mapsto\bar z/z$；
\item[2)] 对每个 $h\in X$，群 $\{g\in G(\C)\mid h(i)gh(i)^{-1}=g\}$ 模其中心是紧的；
\item[3)] 伴随群 $G^{\mathrm{ad}}$ 没有定义在 $\Q$ 上、且 $h\in X$ 的投影在其上平凡的单因子。
\end{itemize}
\item 例如取 $G=\GL_2$，$X$ 为通常嵌入 $\C^\times\to\GL_2(\R)$ 的共轭类。注意 $X\cong\HH^\pm$。
\item 一般地，对 $h\in X$，$h$ 在 $G(\R)$ 中的稳定化子 $K_h$ 模 $G(\R)$ 的中心是紧的（这由条件 2) 推出），且 $X\cong G(\R)/K_h$（$\cong G(\R)/A_GK_\infty$）。
\end{itemize}

% ===== PDF p.101 =====
\begin{itemize}
\item 对 $K\subseteq G(\Af)$，回忆
\[
X_K:=G(\Q)\backslash\bigl(G(\Af)/K\times X\bigr).
\]
另一种记号：$Sh_K(G,X)$，即附着于 $(G,X)$ 的级 $K$ 的 \textbf{Shimura 簇}。
\item 我们得到：
\begin{itemize}
\item[(1)] $X$ 是 Hermite 对称域的无交并，参见 \texttt{[Mil05]}、\texttt{[Del79]}。
\item[(2)] 对每个 $h\in X$ 自然地附着抛物子群 $P_h\subseteq G(\C)$（参见 \texttt{[CS17, Section 2.1]}），且 $h\mapsto P_h$ 定义了开嵌入
\[
\pi\colon X\cong G(\R)/K_h\longrightarrow\mathcal{Fl}\cong G(\C)/P_h,
\]
即把 $X$ 嵌入一个旗簇（这推广了嵌入 $\HH^\pm\to\Pone(\C)$）。\znoteplain{原文脚注（16）：把 $X$ 视为「（极化）Hodge 结构的模空间」可以解释这个嵌入，参见 \texttt{[Del79]}。}
特别地，$X$ 带有自然的复结构。更强的是（Baily–Borel）：$X_K$ 自然地是拟射影簇，参见 \texttt{[Lan17, Theorem 2.4.1.]}。
\item[(3)] 旗簇\idx{旗簇} $\mathcal{Fl}$ 上 $G(\C)$-等变向量丛的拉回给出 $G(\Af)\times X$ 上 $G(\Q)$-等变向量丛。经下降，对任意紧开拓 $K\subseteq G(\Af)$ 得到 $X_K$ 上的\textbf{自守向量丛\idx{自守向量丛}} $\mathcal E$，即 $\omega_k$（$k\in\Z$）的类比，参见 \texttt{[Har88]}。注意 $\mathcal{Fl}$ 上 $G(\C)$-等变向量丛等价于 $P_h$ 的（代数）表示。
\item[(4)、(5)] 若 $\dim_\C(X_K)>1$，则确实存在紧化 $X_K$，但它们不再是典范的。不过，可以在 $H^*(X_K,\mathcal E)$ 的上同调中定义子空间，以推广 $k\in\Z$ 的 $M_k$，参见 \texttt{[Har88]}。
\item[(6)] 空间 $\widehat X$（= 复流形的逆向极限）再次作为概形 $\widehat X\to\Spec(E)$ 的 $\C$-点出现，该概形即\textbf{典范模型\idx{典范模型}}，定义在某个数域（「Shimura 数据的反射域」）上，参见 \texttt{[Lan17, Theorem 2.4.3.]}。
\item[(8)、(9)、(10)] 有了典范模型，就可以尝试（有时会成功）得到第 (8)、(9)、(10) 条的类比，并最终把（某些）自守表示与（各种变体的）Galois 表示联系起来，参见 \texttt{[Har88]}、\texttt{[BRb]}。不用说，一切都（紧化、Eichler–Shimura 同构/关系、整模型、mod $p$ 几何、内上同调与相交上同调……）要复杂得多！\znote{原文此处作 “obtain analogs of points (18),(19),(20)”，与本条开头的 “(8),(9),(10)” 不一致，按上下文此处应为 (8)、(9)、(10)，译文据此处理。又，原文 “By descent one one obtains” 衍一 “one” 字，译文按 “one obtains” 处理。}
\end{itemize}
\end{itemize}
